\section{Step 5: Turn the no-incoming condition into a pole equation}
\label{sec:wkb-resonance}
% BEGIN CURATED SUBJECT INDEX
\index{boundary condition!incoming and outgoing}
% END CURATED SUBJECT INDEX

\subsection{Outgoing conditions as subdominance after continuation}
\label{subsec:wkb-outgoing}
% BEGIN CURATED SUBJECT INDEX
\index{spectrum!essential}
\index{boundary condition!incoming and outgoing}
\index{exact WKB}
\index{complex scaling}
% END CURATED SUBJECT INDEX

On the real axis a resonance wave grows at both spatial ends.  In the complex
$x$ plane, the same exponential is decaying in appropriate sectors.  Exact
WKB analysis selects a Borel-summed solution that is subdominant in one such
sector, transports it across the Stokes graph, and requires it to be
subdominant in the sector representing the other outgoing end.  Complex
scaling implements a fixed contour through these sectors; exact WKB records
the sector changes algebraically.

Suppose $V(x)\to0$ at both ends and define $k$ by
Eq.~\eqref{eq:energy-momentum}.  Asymptotically,
$\vsub{S}{odd}\to\rmi k$ or $-\rmi k$.  Rotating the coordinate by
$x=\rme^{\rmi\theta}y$ changes which exponential is subdominant according to
Eq.~\eqref{eq:scaled-outgoing}.  The exposure condition of complex scaling is
therefore a statement that the chosen contour lies inside the exact-WKB
subdominant sectors.  From this viewpoint, rotation of the essential spectrum
and deformation of the Stokes graph are complementary descriptions of the
same analyticity
\cite{Morikawa:2025xjq,Morikawa:2025vvs}.

\subsection{The inverted Rosen--Morse barrier}
\label{subsec:rosen-morse}
% BEGIN CURATED SUBJECT INDEX
\index{connection coefficient}
\index{boundary condition!incoming and outgoing}
\index{Stokes curve}
\index{turning point}
\index{Rosen--Morse potential}
% END CURATED SUBJECT INDEX

The inverted Rosen--Morse, or P\"oschl--Teller barrier, is
\begin{equation}
  V(x)=\frac{U_0}{\cosh^2(\beta x)},
  \qquad U_0>0,
  \qquad \beta>0 .
  \label{eq:rosen-morse-potential}
\end{equation}
It is short ranged and exactly solvable, yet its analytic continuation has an
infinite periodic structure of poles and turning points.  It therefore tests
both scattering and exact-WKB descriptions without reducing the problem to a
polynomial potential
\cite{PhysRev.42.210,Zhao:1995,Landau:1991wop,Ryaboy:1993}.

Introduce
\begin{align}
  \xi&=\tanh(\beta x),
  &E&=\frac{\hbar^2k^2}{2m},
  \notag\\
  s(s+1)&=-\frac{2mU_0}{\beta^2\hbar^2} .
  \label{eq:rosen-morse-parameters}
\end{align}
The Schr\"odinger equation becomes an associated Legendre equation.  A
solution outgoing at $x\to+\infty$ may be written in terms of a Gauss
hypergeometric function.  Its incident coefficient at $x\to-\infty$ is,
up to a nonzero convention-dependent phase,
\begin{equation}
  \vsub{A}{in}(k)
  \propto
  \frac{\Gamma(1-\rmi k/\beta)
  \Gamma(-\rmi k/\beta)}
  {\Gamma(-\rmi k/\beta-s)
  \Gamma(-\rmi k/\beta+s+1)} .
  \label{eq:rosen-morse-incoming}
\end{equation}
The resonance condition $\vsub{A}{in}(k)=0$ follows from a pole of either
Gamma function in the denominator.

For a barrier satisfying
\begin{equation}
  \frac{8mU_0}{\beta^2\hbar^2}>1,
  \label{eq:strong-barrier}
\end{equation}
define the positive dimensionless number
\begin{equation}
  \kappa
  =\sqrt{\frac{8mU_0}{\beta^2\hbar^2}-1} .
  \label{eq:kappa-rosen-morse}
\end{equation}
The resonance momenta in the fourth quadrant are
\begin{equation}
  k_n^{\mathrm R}
  =\frac{\beta}{2}
  \left[\kappa-\rmi(2n+1)\right],
  \qquad n=0,1,2,\ldots ,
  \label{eq:rosen-morse-k}
\end{equation}
and the corresponding energies are
\begin{equation}
  E_n^{\mathrm R}
  =\frac{\hbar^2\beta^2}{8m}
  \left[\kappa-\rmi(2n+1)\right]^2 .
  \label{eq:rosen-morse-energy}
\end{equation}
The antiresonant partners obey
$\vsup{k_n}{R}=-(\vsup{k_n}{AR})^*$ for a real potential.

In exact-WKB language, the turning points satisfy
\begin{equation}
  \cosh^2(\beta x)=\frac{U_0}{E} .
  \label{eq:rosen-morse-turning-points}
\end{equation}
They repeat under imaginary shifts because $\cosh z$ is periodic up to sign.
The spectral curve also has double poles at zeros of $\cosh(\beta x)$.  A
connection path crosses the relevant Stokes curves, changes sheets at branch
cuts, and returns through infinity.  The exact connection coefficient can be
identified with Eq.~\eqref{eq:rosen-morse-incoming}; hence the resummed WKB
condition reproduces Eqs.~\eqref{eq:rosen-morse-k} and
\eqref{eq:rosen-morse-energy} exactly
\cite{Morikawa:2025grx,Morikawa:2025xjq}.

The lesson is broader than solvability.  Barrier resonance does not require a
real-axis potential well.  In complex configuration space, turning points and
subdominant sectors can trap analytic continuation around a nontrivial cycle.
The pole measures that global complex geometry.  This statement should not be
misread as a sufficient criterion based only on visual minima of $\re V(z)$;
the invariant criterion is the vanishing connection coefficient.

\subsection{Transmission and reflection from connection matrices}
\label{subsec:wkb-transmission}
% BEGIN CURATED SUBJECT INDEX
\index{Wronskian}
\index{scattering matrix}
\index{connection coefficient}
\index{boundary condition!incoming and outgoing}
\index{resonance!residue}
\index{exact WKB}
% END CURATED SUBJECT INDEX

Choose an incoming wave of unit amplitude from the left.  The asymptotic
solution is
\begin{align}
  \psi(x)&\sim
  \rme^{\rmi kx}+R(k)\rme^{-\rmi kx},
  &&x\to-\infty,
  \notag\\
  \psi(x)&\sim T(k)\rme^{\rmi kx},
  &&x\to+\infty .
  \label{eq:RT-boundary}
\end{align}
Let $M(E,\hbar)$ be the exact connection matrix between asymptotic WKB bases.
With a consistent normalization, $R$ and $T$ are ratios of its entries.  The
resonance condition is the zero of the entry multiplying the incoming wave
when the right solution is chosen outgoing.  On the real axis, Wronskian
conservation gives
\begin{equation}
  |R(E)|^2+|T(E)|^2=1
  \label{eq:unitarity-RT}
\end{equation}
for a real one-channel potential with equal asymptotic velocities.  Exact WKB
therefore reconstructs both the continued poles and the unitary physical
scattering matrix from the same connection data
\cite{Morikawa:2025vvs}.

Near a simple pole, the connection coefficient has the local expansion
\begin{equation}
  \Cincoming(E)
  =\Cincoming'(\resonantE)
  (E-\resonantE)+\order((E-\resonantE)^2) .
  \label{eq:connection-local-pole}
\end{equation}
The derivative $\Cincoming'(\resonantE)$ determines the residue.  This
is the exact-WKB route to excitation strength, and not only to the pole
location.

\subsection{Complex scaling as a change of Stokes sector}
\label{subsec:abc-wkb}
% BEGIN CURATED SUBJECT INDEX
\index{unbounded operator}
\index{spectrum!continuous}
\index{spectrum!essential}
\index{resolvent}
\index{exact WKB}
\index{complex scaling}
% END CURATED SUBJECT INDEX

For an asymptotically free problem, the continuous spectrum corresponds to
oscillatory WKB solutions on the real axis.  Under $x=\rme^{\rmi\theta}y$, the
classical action is continued to a ray on which one solution decays and the
other grows.  The energy at which neither is exponentially preferred lies on
the rotated ray $\arg E=-2\theta$.  Thus the continuum rotation in
Eq.~\eqref{eq:continuum-rotation} follows directly from the asymptotic WKB
exponent.  A pole is $\theta$ independent because its global connection
condition is invariant while the contour remains in the same homology class
and crosses no singularity.

This gives an alternative geometric reading of the ABC structure.  It does
not replace the operator theorem: exact WKB controls differential-equation
solutions and their analytic continuation, while the theorem controls
operator domains, essential spectrum, and resolvent matrix elements.  Their
agreement is powerful precisely because the two analyses have different
failure modes.

\subsection{Continuum states in exact-WKB language}
\label{subsec:wkb-continuum}
% BEGIN CURATED SUBJECT INDEX
\index{spectrum!point}
\index{resolvent}
\index{boundary condition!incoming and outgoing}
\index{resonance!residue}
\index{Borel summation}
\index{exact WKB}
\index{complex scaling}
% END CURATED SUBJECT INDEX

At real $E>0$, a continuum state is specified by a boundary value
$E\pm\rmi0$ and by incoming channel amplitudes.  In exact WKB, these choices
select a lateral Borel sum and a side of the Stokes graph.  The continuum is
therefore not obtained by solving the resonance condition for a dense set of
energies; it is a family of connection problems parametrized by real $E$.

After complex scaling, the continuum becomes a line of generalized
eigenvalues $E=\rme^{-2\rmi\theta}\hbar^2q^2/(2m)$ with $q\ge0$.  A finite
basis discretizes $q$, while exact WKB retains it as a continuous parameter.
The equivalence between the two descriptions is clearest at the resolvent
level: both deform the same continuum contour while keeping pole residues
inside the deformation fixed.
